\documentclass[11pt,a4paper]{article}
\usepackage[utf8]{inputenc}
\usepackage[turkish]{babel}
\usepackage[T1]{fontenc}
\usepackage{geometry}
\usepackage{longtable}
\usepackage{array}
\usepackage{hyperref}
\usepackage{xcolor}
\geometry{left=2.2cm,right=2.2cm,top=2.2cm,bottom=2.2cm}

\title{University1652-Fork\\Dosya Özeti (LaTeX)}
\author{Otomatik Proje Envanteri}
\date{\today}

\begin{document}
\maketitle

\section*{Amaç}
Bu doküman, \texttt{University1652-Fork} çalışma alanındaki dosya ve klasörlerin kısa görev özetini verir.
Odak: dosya bazında ``ne işe yarıyor?'' sorusuna hızlı yanıt.

\section*{Kök Dizin Dosyaları}
\begin{longtable}{|p{0.34\textwidth}|p{0.60\textwidth}|}
\hline
\textbf{Dosya} & \textbf{Özet} \\
\hline
\endfirsthead
\hline
\textbf{Dosya} & \textbf{Özet} \\
\hline
\endhead

autoaugment.py & Veri artırma (AutoAugment benzeri) politikaları; eğitimde çeşitlilik sağlar. \\
\hline
circle_loss.py & Metric learning için Circle Loss implementasyonu. \\
\hline
colab_full_pipeline.py & Colab ortamında uçtan uca eğitim/değerlendirme akışını çalıştırma betiği. \\
\hline
COLAB_SATELLITE_DRONE.ipynb & Colab notebook; deney, görselleştirme ve hızlı test için etkileşimli akış. \\
\hline
compare_gradcam.py & Farklı model/konfigürasyonların Grad-CAM çıktılarını karşılaştırma. \\
\hline
compare_models.py & Modellerin metrik ve/veya tahmin bazlı kıyaslaması. \\
\hline
dataset_rgbd.py & RGB+Depth (4 kanal) veri yükleme/dönüştürme mantığı. \\
\hline
demo.py & Standart çözünürlükte örnek çıkarım/demo çalıştırma. \\
\hline
demo_4K.py & 4K görseller için demo/çıkarım akışı. \\
\hline
evaluate_gpu.py & GPU üzerinde değerlendirme (retrieval/accuracy) betiği. \\
\hline
find_rgb_improvements.py & RGB ve RGBD sonuçlarını karşılaştırıp iyileşen örnekleri bulma ve görselleştirme. \\
\hline
folder.py & Veri klasörlerini organize etme/okuma yardımcıları. \\
\hline
gallery_name.txt & Gallery görsel isim listesi. \\
\hline
gradcam_visualization.py & Grad-CAM üretimi, ön işleme ve ısı haritası bindirme yardımcıları. \\
\hline
LICENSE & Lisans metni. \\
\hline
model.py & Temel model tanımları (ResNet/VGG tabanlı, GeM dahil, iki/üç görünüm ağları). \\
\hline
model_rgbd.py & RGBD modeli (4 kanallı giriş), RGB ve RGBD iki görünüm mimarisi. \\
\hline
new_name_list.txt & Yeniden adlandırılmış/işlenmiş dosya isim listesi. \\
\hline
prepare_cvusa.py & CVUSA veri hazırlık/ön işleme betiği. \\
\hline
prepare_limit_view.py & Görüş açısı/sahne kısıtlı veri hazırlama adımları. \\
\hline
query_name.txt & Query görsel isim listesi. \\
\hline
random_erasing.py & Random Erasing veri artırma implementasyonu. \\
\hline
README.md & Projenin genel açıklaması ve kullanım özeti. \\
\hline
README_GRADCAM.md & Grad-CAM kullanım/dokümantasyon notları. \\
\hline
README_RGBD.md & RGBD eğitim/değerlendirme açıklamaları. \\
\hline
Request.md & İstek/notlar veya görev tanımı dokümanı. \\
\hline
requirement.txt & Python bağımlılık listesi (requirements). \\
\hline
show_data.py & Veri örneklerini görselleştirme/inceleme aracı. \\
\hline
technical_documentation.tex & Ayrıntılı teknik dokümantasyon (teori + mimari + pipeline). \\
\hline
test.py & Genel test/evaluation giriş betiği. \\
\hline
test_160k.py & 160k veri/konfigürasyon için test betiği. \\
\hline
test_4K.py & 4K veri üstünde test/değerlendirme. \\
\hline
test_cvusa.py & CVUSA özel test scripti. \\
\hline
test_ImageNet.py & ImageNet benzeri akışta test/deneme betiği. \\
\hline
test_Place365.py & Place365 benzeri akışta test/deneme betiği. \\
\hline
test_train.py & Eğitim doğrulama/test kontrol scripti. \\
\hline
train.py & Ana eğitim scripti. \\
\hline
train_1_sample.py & Çok küçük örnekle (1 örnek) debug eğitim akışı. \\
\hline
train_3_sample.py & 3 örnekle hızlı debug eğitim akışı. \\
\hline
train_9_sample.py & 9 örnekle hızlı debug eğitim akışı. \\
\hline
train_18_sample.py & 18 örnekle hızlı debug eğitim akışı. \\
\hline
train_27_sample.py & 27 örnekle hızlı debug eğitim akışı. \\
\hline
train_baseline_rgb.sh & RGB baseline eğitim komutlarını içeren shell script. \\
\hline
train_cvusa.py & CVUSA veri seti için eğitim scripti. \\
\hline
train_no_street.py & Sokak görünümü hariç/sınırlı veriyle eğitim deneyi. \\
\hline
train_rgbd.sh & RGBD eğitim komutlarını içeren shell script. \\
\hline
TRAINING_GUIDE.md & Eğitim adımları ve pratik öneriler. \\
\hline
utils.py & Ortak yardımcı fonksiyonlar (I/O, metrik, vb.). \\
\hline
\end{longtable}

\section*{Alt Klasörler ve İçerikleri}
\subsection*{docs/}
\begin{longtable}{|p{0.34\textwidth}|p{0.60\textwidth}|}
\hline
\textbf{Dosya/Klasör} & \textbf{Özet} \\
\hline
\endfirsthead
\hline
\textbf{Dosya/Klasör} & \textbf{Özet} \\
\hline
\endhead

docs/index.html & Dokümantasyon ana HTML sayfası. \\
\hline
docs/index_files/ & HTML için yardımcı statik dosyalar. \\
\hline
docs/index_files/bibtex_cvpr2019.txt & CVPR 2019 bibtex referansı. \\
\hline
\end{longtable}

\subsection*{GPU-Re-Ranking/}
\begin{longtable}{|p{0.34\textwidth}|p{0.60\textwidth}|}
\hline
\textbf{Dosya/Klasör} & \textbf{Özet} \\
\hline
\endfirsthead
\hline
\textbf{Dosya/Klasör} & \textbf{Özet} \\
\hline
\endhead

GPU-Re-Ranking/evaluate_rerank_gpu.py & GPU tabanlı re-ranking değerlendirmesi. \\
\hline
GPU-Re-Ranking/gnn_reranking.py & GNN tabanlı re-ranking algoritması. \\
\hline
GPU-Re-Ranking/utils.py & Re-ranking yardımcı fonksiyonları. \\
\hline
GPU-Re-Ranking/README.md & Re-ranking modülü kullanımı. \\
\hline
GPU-Re-Ranking/extension/make.sh & CUDA/C++ extension derleme scripti. \\
\hline
GPU-Re-Ranking/extension/adjacency_matrix/build_adjacency_matrix.cpp & Adjacency matrix extension (C++ köprüsü). \\
\hline
GPU-Re-Ranking/extension/adjacency_matrix/build_adjacency_matrix_kernel.cu & Adjacency matrix CUDA kernel. \\
\hline
GPU-Re-Ranking/extension/adjacency_matrix/setup.py & Adjacency extension kurulum dosyası. \\
\hline
GPU-Re-Ranking/extension/propagation/gnn_propagate.cpp & GNN propagate extension (C++ köprüsü). \\
\hline
GPU-Re-Ranking/extension/propagation/gnn_propagate_kernel.cu & GNN propagation CUDA kernel. \\
\hline
GPU-Re-Ranking/extension/propagation/setup.py & Propagation extension kurulum dosyası. \\
\hline
\end{longtable}

\subsection*{SIFT/}
\begin{longtable}{|p{0.34\textwidth}|p{0.60\textwidth}|}
\hline
\textbf{Dosya} & \textbf{Özet} \\
\hline
\endfirsthead
\hline
\textbf{Dosya} & \textbf{Özet} \\
\hline
\endhead

SIFT/demo_sift.m & MATLAB SIFT demo dosyası. \\
\hline
SIFT/README.md & SIFT modülü açıklaması. \\
\hline
\end{longtable}

\subsection*{State-of-the-art/}
\begin{longtable}{|p{0.34\textwidth}|p{0.60\textwidth}|}
\hline
\textbf{Dosya} & \textbf{Özet} \\
\hline
\endfirsthead
\hline
\textbf{Dosya} & \textbf{Özet} \\
\hline
\endhead

State-of-the-art/README.md & Literatür/karşılaştırmalı sonuç notları. \\
\hline
\end{longtable}

\subsection*{tool/ ve tutorial/}
\begin{longtable}{|p{0.34\textwidth}|p{0.60\textwidth}|}
\hline
\textbf{Dosya} & \textbf{Özet} \\
\hline
\endfirsthead
\hline
\textbf{Dosya} & \textbf{Özet} \\
\hline
\endhead

tool/clear_model.py & Eski/ağır model dosyalarını temizleme aracı. \\
\hline
tutorial/README.md & Hızlı başlangıç/eğitim öğretici notları. \\
\hline
\end{longtable}

\subsection*{Varlık ve Veri Klasörleri}
\begin{longtable}{|p{0.34\textwidth}|p{0.60\textwidth}|}
\hline
\textbf{Klasör} & \textbf{Özet} \\
\hline
\endfirsthead
\hline
\textbf{Klasör} & \textbf{Özet} \\
\hline
\endhead

image_4K/00..11 & 4K örnek görüntü klasörleri (demo/test girdileri). \\
\hline
model/use_rgbd/net_059.pth & RGBD eğitim checkpoint ağırlığı. \\
\hline
model/use_rgbd/net_079.pth & RGBD eğitim checkpoint ağırlığı. \\
\hline
model/use_rgbd/net_099.pth & RGBD eğitim checkpoint ağırlığı. \\
\hline
__pycache__/ & Python bytecode önbellek dosyaları. \\
\hline
\end{longtable}

\section*{Notlar}
\begin{itemize}
    \item Bu özet, çalışma alanında görünen dosya ağacına göre hazırlanmıştır.
    \item Bazı scriptler benzer amaçlı varyasyonlardır (ör. farklı veri seti, örnek sayısı veya çözünürlük).
    \item En güncel ve ayrıntılı teorik bilgi için \texttt{technical\_documentation.tex} dosyasına bakın.
\end{itemize}

\end{document}
